% !TEX root = ../jlt.tex
\section{Measurements}
\label{app:measurements}

This appendix measures Section~\ref{sec:binding}'s claim that the compiler vectorises the test a law writes.
A vector register is a direct product of lanes, and its operations act pointwise.
A computation fits it when two conditions hold: no position depends on another, and the values are contiguous and small, under operations the unit has.
Either can fail alone, as Table~\ref{tab:batch} shows.
Listing~\ref{lst:sites}'s table is Brzozowski's automaton~\cite{brzozowski1964derivatives}: for one word the matcher of Knuth, Morris and Pratt~\cite{knuth1977fast}, for many Aho and Corasick's~\cite{aho1975efficient}.
It is generic over the alphabet, and it fails the first condition, as each state needs the last.
\cppinline{fixed_by(~rev)} meets it: the ranks embed the alphabet with its hole in a byte and carry the reflection to $3 - r$, so a pair is a sum compared with three.
\cppinline{count} has an overload keyed on the section's type.
A block goes through in stages, ranks, a pass per pair and a sum, each a fold small enough to be seen as a loop.
It is an exhibit for this compiler and these flags.
A specialisation is keyed on a law or a data shape, never on a particular problem or compiler; it guarantees the answer, while vectorising stays the compiler's choice.
Generating such code from algebra is old~\cite{frigo2005fftw,puschel2005spiral}, as is separating an algorithm from its schedule~\cite{ragankelley2013halide} and matching many patterns in lanes~\cite{wang2019hyperscan}.
What is shown here is that standard \Cpp{23} gets there from laws its types already carry.
The times, the disassembly and what else was tried are in the repository's \texttt{src/bench}, and \texttt{make bench} reproduces them.
The figures are to one significant digit, since benchmarks are fickle and the orders of magnitude are the result.
Past the cache, the law's scan is bound by the core, not by memory: it runs at a tenth of a plain read's rate.

\begin{table}[htbp]
\centering
\small
\begin{tabular}{@{}lccr@{}}
\toprule
Counting the sites of \texttt{pali6} & independent & batchable & ns a letter \\
\midrule
each of the 64 sites' windows tested & yes & no & 20 \\
the 64 sites' table (277 states) & no & yes & 2 \\
\quad the same table, hand-written, holes as the library & no & yes & 2 \\
the law, letters as they come, four compared at run time & yes & no & 10 \\
the law, sets of bases in bytes (IUPAC's code) & yes & yes, with costly operations & 4 \\
\textbf{the library's \texttt{count(text >> pali6)}, by the law} & yes & yes & \textbf{0.4} \\
\quad \texttt{pynzli}, by the law, from bytes & yes & yes & 0.4 \\
\texttt{re}, 64 alternatives in a lookahead & & & 300 \\
\bottomrule
\end{tabular}
\caption{The two conditions, one failing at a time.
Each row is named as \texttt{make bench} prints it.
The figures are one recorded run, in full under \emph{One run in full} in \texttt{src/bench/README.md}: one core of a shared AMD EPYC 7763, Clang~23 at \texttt{-O2}, baseline SSE2, $10^8$ letters, to one significant digit.
Past the cache, on 3.1 billion letters: the law 2\,GB/s, a plain read of the same memory 20\,GB/s.
A text drawn and not stored, by the table: peak memory flat from $10^7$ to $10^{10}$ letters; the law's own storage a fixed block.}
\label{tab:batch}
\end{table}
